\documentclass[11pt,a4paper]{article}
% ==============================================================
%  Préambule commun aux fiches -- Advanced Probabilities (M1)
%  Usage, dans chaque fiche :
%     \documentclass[11pt,a4paper]{article}
%     % ==============================================================
%  Préambule commun aux fiches -- Advanced Probabilities (M1)
%  Usage, dans chaque fiche :
%     \documentclass[11pt,a4paper]{article}
%     % ==============================================================
%  Préambule commun aux fiches -- Advanced Probabilities (M1)
%  Usage, dans chaque fiche :
%     \documentclass[11pt,a4paper]{article}
%     \input{../preambule.tex}
%     \fiche{C5}{Vecteurs gaussiens}{Gaussian vectors}
% ==============================================================
\usepackage[T1]{fontenc}
\usepackage[utf8]{inputenc}
\usepackage[french]{babel}
\usepackage{lmodern}
\usepackage[a4paper,margin=1.9cm,top=1.7cm,bottom=1.9cm]{geometry}
\usepackage{amsmath,amssymb,amsthm,mathtools}
\usepackage[inline]{enumitem}
\usepackage{booktabs,array,tabularx}
\usepackage{xcolor}
\usepackage[most]{tcolorbox}
\usepackage{fancyhdr}
\usepackage{hyperref}
\hypersetup{colorlinks=true,linkcolor=blue!50!black,urlcolor=blue!50!black}

\setlist{itemsep=1pt,topsep=3pt,parsep=0pt}
\setlength{\parindent}{0pt}
\setlength{\parskip}{3pt}
\emergencystretch=3em

% ---------- Identification de la fiche -------------------------
\newcommand{\ficheid}{}
\newcommand{\fichetitre}{}
\newcommand{\fiche}[3]{%
  \renewcommand{\ficheid}{#1}\renewcommand{\fichetitre}{#2}%
  \pagestyle{fancy}\fancyhf{}%
  \fancyhead[L]{\small\textsf{Advanced Probabilities -- Fiche #1}}%
  \fancyhead[R]{\small\textsf{#2}}%
  \fancyfoot[C]{\small\thepage}%
  \begin{center}
    {\LARGE\bfseries Fiche #1 -- #2}\par\smallskip
    {\large\itshape (#3)}\par\smallskip
    {\small Cours de V.\,M.~Lambert, M1 -- semestre 1}
  \end{center}\vspace{-2pt}\hrule\vspace{6pt}}

% ---------- Environnements numérotés (un seul compteur) --------
\theoremstyle{definition}
\newtheorem{definition}{Définition}
\renewcommand{\thedefinition}{\ficheid.\arabic{definition}}
\newtheorem{theoreme}[definition]{Théorème}
\newtheorem{proposition}[definition]{Proposition}
\newtheorem{lemme}[definition]{Lemme}
\newtheorem{corollaire}[definition]{Corollaire}
\newtheorem{exemple}[definition]{Exemple}
\theoremstyle{remark}
\newtheorem{remarque}[definition]{Remarque}

% Résultat hors cours (standard, signalé comme tel)
\theoremstyle{definition}
\newtheorem{horscours}[definition]{Hors cours}

% Idée de preuve
\newenvironment{idee}{\par\small\textit{Idée de preuve.}~}{\par}

% ---------- Encadrés -----------------------------------------
\tcbset{fichebox/.style={enhanced,breakable,boxrule=0.6pt,arc=2pt,
  left=5pt,right=5pt,top=3pt,bottom=3pt,fonttitle=\bfseries}}
\newtcolorbox{piege}[1][]{fichebox,colback=red!4,colframe=red!60!black,
  title={Pièges classiques},#1}
\newtcolorbox{liens}[1][]{fichebox,colback=blue!3,colframe=blue!50!black,
  title={Liens entre notions},#1}
\newtcolorbox{methode}[1][Méthode]{fichebox,colback=green!4,colframe=green!40!black,title={#1}}
\newtcolorbox{alerte}[1][]{fichebox,colback=orange!6,colframe=orange!70!black,
  title={Attention à l'énoncé},#1}

% Exercice de TD (fiches de rappels)
\newcounter{exo}
\newenvironment{exo}[1][]{\refstepcounter{exo}\par\medskip
  \noindent{\large\bfseries Exercice~\theexo}\ifx\relax#1\relax\else~{\large\bfseries -- #1}\fi\par\nopagebreak\smallskip}{\par}
\newcommand{\outils}{\textbf{Outils du cours.}~}
\newcommand{\meth}{\textbf{Méthode / astuce.}~}
\newcommand{\indic}{\textbf{Indication.}~}
\newcommand{\rep}{\textbf{Réponse.}~}

% ---------- Notations du cours ---------------------------------
\newcommand{\R}{\mathbb{R}}
\newcommand{\N}{\mathbb{N}}
\newcommand{\Q}{\mathbb{Q}}
\newcommand{\C}{\mathbb{C}}
\newcommand{\PP}{\mathbb{P}}
\newcommand{\QQ}{\mathbb{Q}}
\newcommand{\E}{\mathbb{E}}
\newcommand{\F}{\mathcal{F}}
\newcommand{\G}{\mathcal{G}}
\newcommand{\Hc}{\mathcal{H}}
\newcommand{\Tc}{\mathcal{T}}
\newcommand{\Bc}{\mathcal{B}}
\newcommand{\Cc}{\mathcal{C}}
\newcommand{\Ec}{\mathcal{E}}
\newcommand{\Pc}{\mathcal{P}}
\newcommand{\Nc}{\mathcal{N}}
\newcommand{\LN}{\mathcal{LN}}
\newcommand{\Expo}{\mathcal{E}}
\newcommand{\Poi}{\mathcal{P}}
\newcommand{\Ff}{\mathbb{F}}
\newcommand{\1}{\mathbf{1}}
\newcommand{\dd}{\,\mathrm{d}}
\DeclareMathOperator{\Var}{\mathbb{V}}
\DeclareMathOperator{\Cov}{Cov}
\DeclareMathOperator{\card}{card}
\DeclareMathOperator{\diag}{diag}
\newcommand{\indep}{\perp\!\!\!\perp}
\newcommand{\ps}{\ \PP\text{-p.s.}}
\newcommand{\tendloi}{\xrightarrow[n\to+\infty]{\text{loi}}}
\newcommand{\tendps}{\xrightarrow[n\to+\infty]{\text{p.s.}}}
\newcommand{\tendproba}{\xrightarrow[n\to+\infty]{\PP}}
\newcommand{\tendLp}{\xrightarrow[n\to+\infty]{L^p}}
\newcommand{\tend}{\xrightarrow[n\to+\infty]{}}

%     \fiche{C5}{Vecteurs gaussiens}{Gaussian vectors}
% ==============================================================
\usepackage[T1]{fontenc}
\usepackage[utf8]{inputenc}
\usepackage[french]{babel}
\usepackage{lmodern}
\usepackage[a4paper,margin=1.9cm,top=1.7cm,bottom=1.9cm]{geometry}
\usepackage{amsmath,amssymb,amsthm,mathtools}
\usepackage[inline]{enumitem}
\usepackage{booktabs,array,tabularx}
\usepackage{xcolor}
\usepackage[most]{tcolorbox}
\usepackage{fancyhdr}
\usepackage{hyperref}
\hypersetup{colorlinks=true,linkcolor=blue!50!black,urlcolor=blue!50!black}

\setlist{itemsep=1pt,topsep=3pt,parsep=0pt}
\setlength{\parindent}{0pt}
\setlength{\parskip}{3pt}
\emergencystretch=3em

% ---------- Identification de la fiche -------------------------
\newcommand{\ficheid}{}
\newcommand{\fichetitre}{}
\newcommand{\fiche}[3]{%
  \renewcommand{\ficheid}{#1}\renewcommand{\fichetitre}{#2}%
  \pagestyle{fancy}\fancyhf{}%
  \fancyhead[L]{\small\textsf{Advanced Probabilities -- Fiche #1}}%
  \fancyhead[R]{\small\textsf{#2}}%
  \fancyfoot[C]{\small\thepage}%
  \begin{center}
    {\LARGE\bfseries Fiche #1 -- #2}\par\smallskip
    {\large\itshape (#3)}\par\smallskip
    {\small Cours de V.\,M.~Lambert, M1 -- semestre 1}
  \end{center}\vspace{-2pt}\hrule\vspace{6pt}}

% ---------- Environnements numérotés (un seul compteur) --------
\theoremstyle{definition}
\newtheorem{definition}{Définition}
\renewcommand{\thedefinition}{\ficheid.\arabic{definition}}
\newtheorem{theoreme}[definition]{Théorème}
\newtheorem{proposition}[definition]{Proposition}
\newtheorem{lemme}[definition]{Lemme}
\newtheorem{corollaire}[definition]{Corollaire}
\newtheorem{exemple}[definition]{Exemple}
\theoremstyle{remark}
\newtheorem{remarque}[definition]{Remarque}

% Résultat hors cours (standard, signalé comme tel)
\theoremstyle{definition}
\newtheorem{horscours}[definition]{Hors cours}

% Idée de preuve
\newenvironment{idee}{\par\small\textit{Idée de preuve.}~}{\par}

% ---------- Encadrés -----------------------------------------
\tcbset{fichebox/.style={enhanced,breakable,boxrule=0.6pt,arc=2pt,
  left=5pt,right=5pt,top=3pt,bottom=3pt,fonttitle=\bfseries}}
\newtcolorbox{piege}[1][]{fichebox,colback=red!4,colframe=red!60!black,
  title={Pièges classiques},#1}
\newtcolorbox{liens}[1][]{fichebox,colback=blue!3,colframe=blue!50!black,
  title={Liens entre notions},#1}
\newtcolorbox{methode}[1][Méthode]{fichebox,colback=green!4,colframe=green!40!black,title={#1}}
\newtcolorbox{alerte}[1][]{fichebox,colback=orange!6,colframe=orange!70!black,
  title={Attention à l'énoncé},#1}

% Exercice de TD (fiches de rappels)
\newcounter{exo}
\newenvironment{exo}[1][]{\refstepcounter{exo}\par\medskip
  \noindent{\large\bfseries Exercice~\theexo}\ifx\relax#1\relax\else~{\large\bfseries -- #1}\fi\par\nopagebreak\smallskip}{\par}
\newcommand{\outils}{\textbf{Outils du cours.}~}
\newcommand{\meth}{\textbf{Méthode / astuce.}~}
\newcommand{\indic}{\textbf{Indication.}~}
\newcommand{\rep}{\textbf{Réponse.}~}

% ---------- Notations du cours ---------------------------------
\newcommand{\R}{\mathbb{R}}
\newcommand{\N}{\mathbb{N}}
\newcommand{\Q}{\mathbb{Q}}
\newcommand{\C}{\mathbb{C}}
\newcommand{\PP}{\mathbb{P}}
\newcommand{\QQ}{\mathbb{Q}}
\newcommand{\E}{\mathbb{E}}
\newcommand{\F}{\mathcal{F}}
\newcommand{\G}{\mathcal{G}}
\newcommand{\Hc}{\mathcal{H}}
\newcommand{\Tc}{\mathcal{T}}
\newcommand{\Bc}{\mathcal{B}}
\newcommand{\Cc}{\mathcal{C}}
\newcommand{\Ec}{\mathcal{E}}
\newcommand{\Pc}{\mathcal{P}}
\newcommand{\Nc}{\mathcal{N}}
\newcommand{\LN}{\mathcal{LN}}
\newcommand{\Expo}{\mathcal{E}}
\newcommand{\Poi}{\mathcal{P}}
\newcommand{\Ff}{\mathbb{F}}
\newcommand{\1}{\mathbf{1}}
\newcommand{\dd}{\,\mathrm{d}}
\DeclareMathOperator{\Var}{\mathbb{V}}
\DeclareMathOperator{\Cov}{Cov}
\DeclareMathOperator{\card}{card}
\DeclareMathOperator{\diag}{diag}
\newcommand{\indep}{\perp\!\!\!\perp}
\newcommand{\ps}{\ \PP\text{-p.s.}}
\newcommand{\tendloi}{\xrightarrow[n\to+\infty]{\text{loi}}}
\newcommand{\tendps}{\xrightarrow[n\to+\infty]{\text{p.s.}}}
\newcommand{\tendproba}{\xrightarrow[n\to+\infty]{\PP}}
\newcommand{\tendLp}{\xrightarrow[n\to+\infty]{L^p}}
\newcommand{\tend}{\xrightarrow[n\to+\infty]{}}

%     \fiche{C5}{Vecteurs gaussiens}{Gaussian vectors}
% ==============================================================
\usepackage[T1]{fontenc}
\usepackage[utf8]{inputenc}
\usepackage[french]{babel}
\usepackage{lmodern}
\usepackage[a4paper,margin=1.9cm,top=1.7cm,bottom=1.9cm]{geometry}
\usepackage{amsmath,amssymb,amsthm,mathtools}
\usepackage[inline]{enumitem}
\usepackage{booktabs,array,tabularx}
\usepackage{xcolor}
\usepackage[most]{tcolorbox}
\usepackage{fancyhdr}
\usepackage{hyperref}
\hypersetup{colorlinks=true,linkcolor=blue!50!black,urlcolor=blue!50!black}

\setlist{itemsep=1pt,topsep=3pt,parsep=0pt}
\setlength{\parindent}{0pt}
\setlength{\parskip}{3pt}
\emergencystretch=3em

% ---------- Identification de la fiche -------------------------
\newcommand{\ficheid}{}
\newcommand{\fichetitre}{}
\newcommand{\fiche}[3]{%
  \renewcommand{\ficheid}{#1}\renewcommand{\fichetitre}{#2}%
  \pagestyle{fancy}\fancyhf{}%
  \fancyhead[L]{\small\textsf{Advanced Probabilities -- Fiche #1}}%
  \fancyhead[R]{\small\textsf{#2}}%
  \fancyfoot[C]{\small\thepage}%
  \begin{center}
    {\LARGE\bfseries Fiche #1 -- #2}\par\smallskip
    {\large\itshape (#3)}\par\smallskip
    {\small Cours de V.\,M.~Lambert, M1 -- semestre 1}
  \end{center}\vspace{-2pt}\hrule\vspace{6pt}}

% ---------- Environnements numérotés (un seul compteur) --------
\theoremstyle{definition}
\newtheorem{definition}{Définition}
\renewcommand{\thedefinition}{\ficheid.\arabic{definition}}
\newtheorem{theoreme}[definition]{Théorème}
\newtheorem{proposition}[definition]{Proposition}
\newtheorem{lemme}[definition]{Lemme}
\newtheorem{corollaire}[definition]{Corollaire}
\newtheorem{exemple}[definition]{Exemple}
\theoremstyle{remark}
\newtheorem{remarque}[definition]{Remarque}

% Résultat hors cours (standard, signalé comme tel)
\theoremstyle{definition}
\newtheorem{horscours}[definition]{Hors cours}

% Idée de preuve
\newenvironment{idee}{\par\small\textit{Idée de preuve.}~}{\par}

% ---------- Encadrés -----------------------------------------
\tcbset{fichebox/.style={enhanced,breakable,boxrule=0.6pt,arc=2pt,
  left=5pt,right=5pt,top=3pt,bottom=3pt,fonttitle=\bfseries}}
\newtcolorbox{piege}[1][]{fichebox,colback=red!4,colframe=red!60!black,
  title={Pièges classiques},#1}
\newtcolorbox{liens}[1][]{fichebox,colback=blue!3,colframe=blue!50!black,
  title={Liens entre notions},#1}
\newtcolorbox{methode}[1][Méthode]{fichebox,colback=green!4,colframe=green!40!black,title={#1}}
\newtcolorbox{alerte}[1][]{fichebox,colback=orange!6,colframe=orange!70!black,
  title={Attention à l'énoncé},#1}

% Exercice de TD (fiches de rappels)
\newcounter{exo}
\newenvironment{exo}[1][]{\refstepcounter{exo}\par\medskip
  \noindent{\large\bfseries Exercice~\theexo}\ifx\relax#1\relax\else~{\large\bfseries -- #1}\fi\par\nopagebreak\smallskip}{\par}
\newcommand{\outils}{\textbf{Outils du cours.}~}
\newcommand{\meth}{\textbf{Méthode / astuce.}~}
\newcommand{\indic}{\textbf{Indication.}~}
\newcommand{\rep}{\textbf{Réponse.}~}

% ---------- Notations du cours ---------------------------------
\newcommand{\R}{\mathbb{R}}
\newcommand{\N}{\mathbb{N}}
\newcommand{\Q}{\mathbb{Q}}
\newcommand{\C}{\mathbb{C}}
\newcommand{\PP}{\mathbb{P}}
\newcommand{\QQ}{\mathbb{Q}}
\newcommand{\E}{\mathbb{E}}
\newcommand{\F}{\mathcal{F}}
\newcommand{\G}{\mathcal{G}}
\newcommand{\Hc}{\mathcal{H}}
\newcommand{\Tc}{\mathcal{T}}
\newcommand{\Bc}{\mathcal{B}}
\newcommand{\Cc}{\mathcal{C}}
\newcommand{\Ec}{\mathcal{E}}
\newcommand{\Pc}{\mathcal{P}}
\newcommand{\Nc}{\mathcal{N}}
\newcommand{\LN}{\mathcal{LN}}
\newcommand{\Expo}{\mathcal{E}}
\newcommand{\Poi}{\mathcal{P}}
\newcommand{\Ff}{\mathbb{F}}
\newcommand{\1}{\mathbf{1}}
\newcommand{\dd}{\,\mathrm{d}}
\DeclareMathOperator{\Var}{\mathbb{V}}
\DeclareMathOperator{\Cov}{Cov}
\DeclareMathOperator{\card}{card}
\DeclareMathOperator{\diag}{diag}
\newcommand{\indep}{\perp\!\!\!\perp}
\newcommand{\ps}{\ \PP\text{-p.s.}}
\newcommand{\tendloi}{\xrightarrow[n\to+\infty]{\text{loi}}}
\newcommand{\tendps}{\xrightarrow[n\to+\infty]{\text{p.s.}}}
\newcommand{\tendproba}{\xrightarrow[n\to+\infty]{\PP}}
\newcommand{\tendLp}{\xrightarrow[n\to+\infty]{L^p}}
\newcommand{\tend}{\xrightarrow[n\to+\infty]{}}

\begin{document}
\fiche{C1}{Espaces probabilisés}{Sigma-algebra, Probability measure, Conditional probability}

\textit{Source : slides, chapitres 1 à 3. Notation : les slides notent la tribu $\Tc$, le poly
d'espérance conditionnelle la note $\F$ ; on écrit ici $(\Omega,\F,\PP)$, les deux sont équivalents.}

\section*{1. Tribus (sigma-algebras)}

\begin{definition}[Tribu, \emph{sigma-algebra}]\label{C1:tribu}
Soit $\Omega$ un ensemble. $\F\subset\Pc(\Omega)$ est une \emph{tribu} sur $\Omega$ si
\begin{enumerate}[label=(\roman*)]
  \item $\emptyset\in\F$ ;
  \item $A\in\F \Rightarrow A^c:=\{x\in\Omega,\ x\notin A\}\in\F$ (stabilité par complémentaire) ;
  \item $(A_n)_{n\in\N}\subset\F \Rightarrow \bigcup_{n\in\N}A_n\in\F$ (stabilité par union \emph{dénombrable}).
\end{enumerate}
Les éléments de $\F$ sont les \emph{événements} (\emph{events}). $(\Omega,\F)$ est un \emph{espace mesurable}.
\end{definition}

\begin{remarque}[conséquences immédiates, standard]
$\Omega=\emptyset^c\in\F$ ; $\F$ est stable par intersection dénombrable
($\bigcap A_n=(\bigcup A_n^c)^c$), par différence $A\setminus B=A\cap B^c$ et par union finie.
\end{remarque}

\begin{exemple}[tribus de référence]\label{C1:exemples}\leavevmode
\begin{enumerate}
  \item Tribu triviale $\{\emptyset,\Omega\}$ (aucune information) ;
  \item tribu de toutes les parties $\Pc(\Omega)$ (information complète) ;
  \item tribu associée à une partition $(A_i)_{i\in I}$ de $\Omega$ ($\Omega=\bigsqcup_{i\in I}A_i$) :
  $\bigl\{\bigcup_{j\in J}A_j,\ J\subset I\bigr\}$.
\end{enumerate}
\end{exemple}

\begin{definition}[Tribu engendrée, \emph{sigma-algebra generated by} $\Cc$]\label{C1:engendree}
Pour tout $\Cc\subset\Pc(\Omega)$, il existe une plus petite tribu (au sens de l'inclusion) contenant $\Cc$ ;
on la note $\sigma(\Cc)$.
\end{definition}
\begin{idee}
Une intersection quelconque de tribus est une tribu. On pose $\sigma(\Cc)=\bigcap\{\Tc \text{ tribu},\ \Cc\subset\Tc\}$,
famille non vide car elle contient $\Pc(\Omega)$. (C'est l'objet du TD1, ex.~1.)
\end{idee}

\begin{remarque}[tribu borélienne]
Le cours utilise $\Bc(\R)$ sans la redéfinir. \emph{Hors cours (standard)} : $\Bc(\R)=\sigma(\text{ouverts de }\R)
=\sigma\bigl(\{]-\infty,x],\ x\in\R\}\bigr)$, et $\Bc(\R^d)$ est engendrée par les ouverts de $\R^d$.
\end{remarque}

\begin{definition}[Filtration]\label{C1:filtration}
Soit, pour $t\ge 0$, $\F_t\subset\Pc(\Omega)$ une tribu sur $\Omega$. $(\F_t)_{t\ge0}$ est une
\emph{filtration} si $\F_s\subset\F_t$ dès que $s\le t$. Interprétation : $\F_t$ est l'information disponible à la date $t$.
\end{definition}

\section*{2. Variables aléatoires et processus}

\begin{definition}[Variable aléatoire réelle, \emph{random variable}]\label{C1:va}
Soit $\F$ une tribu sur $\Omega$. Une v.a.\ réelle est une fonction $X:\Omega\to\R$ telle que
\[\forall B\in\Bc(\R),\qquad \{X\in B\}:=X^{-1}(B)=\{\omega\in\Omega,\ X(\omega)\in B\}\in\F,\]
c'est-à-dire que $\{X\in B\}$ est un événement ($X$ est $\F$-mesurable).
\end{definition}

\begin{remarque}[tribu engendrée par $X$, standard]
$\sigma(X):=\{X^{-1}(B),\ B\in\Bc(\R)\}$ est une tribu : c'est la plus petite qui rende $X$ mesurable
(« l'information portée par $X$ »). Le cours l'utilise aux chapitres 5 et 7 ($\E(Y\mid X):=\E(Y\mid\sigma(X))$).
\end{remarque}

\begin{definition}[Processus stochastique adapté]\label{C1:processus}
Soit $(\F_t)_{t\ge0}$ une filtration. Un processus $(X_t)_{t\ge0}$ est une famille de v.a.\ réelles telle que,
pour tout $B\in\Bc(\R)$ et tout $t$, $\{X_t\in B\}\in\F_t$. (Voir la fiche C8.)
\end{definition}

\section*{3. Mesure de probabilité}

\begin{definition}[Probabilité, \emph{probability measure}]\label{C1:proba}
Soit $(\Omega,\F)$ un espace mesurable. $\PP:\F\to\R$ est une probabilité si
\begin{enumerate}[label=(\roman*)]
  \item $\PP(\Omega)=1$ ;
  \item $\PP(A)\ge 0$ pour tout $A\in\F$ ;
  \item ($\sigma$-additivité) si $(A_n)_{n\in\N}\subset\F$ sont \emph{deux à deux disjoints}
  ($i\neq j\Rightarrow A_i\cap A_j=\emptyset$), alors $\PP\bigl(\bigcup_{n}A_n\bigr)=\sum_{n}\PP(A_n)$.
\end{enumerate}
$(\Omega,\F,\PP)$ est un \emph{espace probabilisé} (\emph{probability space}).
\end{definition}

\begin{proposition}[propriétés élémentaires, standard]\label{C1:prop-elem}
Pour $A,B\in\F$ et $(A_n)\subset\F$ :
\begin{enumerate}
  \item $\PP(\emptyset)=0$, \quad $\PP(A^c)=1-\PP(A)$, \quad $A\subset B\Rightarrow \PP(A)\le\PP(B)$ ;
  \item $\PP(A\cup B)=\PP(A)+\PP(B)-\PP(A\cap B)$ (donc $\PP(A\cap B)\ge\PP(A)+\PP(B)-1$) ;
  \item sous-additivité : $\PP\bigl(\bigcup_n A_n\bigr)\le\sum_n\PP(A_n)$ ;
  \item continuité monotone : si $A_n\uparrow A$ alors $\PP(A_n)\uparrow\PP(A)$ ; si $A_n\downarrow A$ alors $\PP(A_n)\downarrow\PP(A)$.
\end{enumerate}
\end{proposition}
\begin{idee}
Tout découle de la $\sigma$-additivité appliquée à des unions disjointes bien choisies :
$A\cup B=A\sqcup(B\setminus A)$, $B=A\sqcup(B\setminus A)$ si $A\subset B$, et pour 4.\ $A=A_0\sqcup\bigsqcup_n(A_{n+1}\setminus A_n)$.
\end{idee}

\section*{4. Probabilité conditionnelle (conditional probability)}

\begin{definition}[Probabilité conditionnelle]\label{C1:condproba}
Dans $(\Omega,\F,\PP)$, pour $A,B\in\F$ avec $\PP(B)>0$ :
\[\PP(A\mid B)=\frac{\PP(A\cap B)}{\PP(B)}.\]
\end{definition}

\begin{horscours}[formules des probabilités totales et de Bayes]\label{C1:bayes}
$A\mapsto\PP(A\mid B)$ est une probabilité sur $(\Omega,\F)$. Si $(B_i)_{i\in I}$ ($I$ fini ou dénombrable)
est une partition de $\Omega$ en événements tous de probabilité $>0$, alors pour tout $A\in\F$
\[\PP(A)=\sum_{i\in I}\PP(A\mid B_i)\PP(B_i),\qquad
\PP(B_j\mid A)=\frac{\PP(A\mid B_j)\PP(B_j)}{\sum_{i}\PP(A\mid B_i)\PP(B_i)}\quad(\text{si }\PP(A)>0).\]
\end{horscours}
\begin{idee}
$A=\bigsqcup_i(A\cap B_i)$ et $\sigma$-additivité ; puis $\PP(B_j\mid A)=\PP(A\cap B_j)/\PP(A)$.
\end{idee}

\begin{liens}
\begin{itemize}
  \item $\{\emptyset,\Omega\}\subset\F_s\subset\F_t\subset\Pc(\Omega)$ : plus la tribu est grosse, plus on a d'information.
  \item Une intersection de tribus est une tribu, \emph{mais} une union de tribus n'en est pas une en général :
  $\Omega=\{1,2,3\}$, $\F_1=\sigma(\{1\})$, $\F_2=\sigma(\{2\})$, et $\{1\}\cup\{2\}=\{1,2\}\notin\F_1\cup\F_2$ (TD1, ex.~1) ; on remplace l'union par $\sigma(\F_1\cup\F_2)$.
  \item Pour une sous-tribu $\G\subset\F$ : « $X$ est $\G$-mesurable » $\iff\sigma(X)\subset\G$. Toute v.a.\ vérifie
  $\sigma(X)\subset\F$ ; être $\G$-mesurable est plus fort (cf.\ fiche C6) : avec $\G=\{\emptyset,\Omega\}$, une v.a.\ non constante est $\F$-mesurable mais pas $\G$-mesurable.
\end{itemize}
\end{liens}

\begin{piege}
\begin{itemize}
  \item Une tribu est stable par union \emph{dénombrable}, pas quelconque (ce n'est pas une topologie !).
  \item La $\sigma$-additivité ne vaut que pour des événements \emph{deux à deux disjoints} ; sinon on n'a que
  la sous-additivité.
  \item $\PP(A\mid B)$ n'est défini que si $\PP(B)>0$, et en général $\PP(A\mid B)\neq\PP(B\mid A)$
  (c'est tout l'objet de la formule de Bayes).
  \item La tribu de partition $\{\bigcup_{j\in J}A_j\}$ coïncide avec $\sigma(\{A_i\})$ quand la partition est
  \emph{dénombrable} (cf.\ TD1, ex.~2) : l'inclusion $\sigma(\{A_i\})\subset\{\bigcup_{j\in J}A_j\}$ est toujours vraie, mais $\{\bigcup_{j\in J}A_j\}\subset\sigma(\{A_i\})$ utilise des unions dénombrables. Pour la partition de $\R$ en singletons, l'égalité est fausse.
  \item Avant d'écrire $\PP(A)$, vérifier que $A\in\F$ ; $\PP$ n'est pas définie sur $\Pc(\Omega)$ en général.
  \item Des nombres « $\PP(A),\PP(B),\PP(A\cap B)$ » arbitraires ne sont pas toujours compatibles :
  il faut $\PP(A)+\PP(B)-\PP(A\cap B)\le 1$ et $\PP(A\cap B)\le\min(\PP(A),\PP(B))$.
\end{itemize}
\end{piege}

\section*{Exemples types}

\begin{exemple}[modèle d'actif à deux périodes, slides ch.~5 et 7]\label{C1:ex-arbre}
$\Omega=\{UU,UD,DU,DD\}$ (hausse/baisse à chaque période), $\F=\Pc(\Omega)$, et pour $0<p<1$ :
$\PP(\{UU\})=p^2$, $\PP(\{UD\})=\PP(\{DU\})=p(1-p)$, $\PP(\{DD\})=(1-p)^2$.
C'est bien une probabilité : les poids sont $\ge0$ et $p^2+2p(1-p)+(1-p)^2=(p+1-p)^2=1$.

\emph{Filtration.} Soit $A_U=\{UU,UD\}$ (« hausse en première période ») et $A_D=A_U^c$.
$\F_0=\{\emptyset,\Omega\}\subset\F_1=\{\emptyset,A_U,A_D,\Omega\}=\sigma(\{A_U\})\subset\F_2=\Pc(\Omega)$ :
$\F_1$ est la tribu de la partition $\{A_U,A_D\}$ et représente l'information à la date 1.

\emph{Conditionnement.} $\PP(A_U)=p^2+p(1-p)=p$, donc
$\PP(\{UU\}\mid A_U)=\dfrac{\PP(\{UU\})}{\PP(A_U)}=\dfrac{p^2}{p}=p$.
\end{exemple}

\begin{exemple}[tribu engendrée par une partie]
Pour $A\subset\Omega$ avec $A\neq\emptyset,\Omega$ : $\sigma(\{A\})=\{\emptyset,A,A^c,\Omega\}$. En effet cette famille
est une tribu (vérifier (i)--(iii)), elle contient $A$, et toute tribu contenant $A$ contient $A^c$, $\emptyset$, $\Omega$.
Pour un dé ($\Omega=\{1,\dots,6\}$) et $A=\{1,3,5\}$, $\sigma(\{A\})=\{\emptyset,\{1,3,5\},\{2,4,6\},\Omega\}$ :
c'est l'information « parité du résultat ».
\end{exemple}

\end{document}
